\section{Layer II: Physical and Systems Analogies}

\subsection{Why Analogy Is Needed, But Must Stay Disciplined}

Biology gives the chapter legitimacy, but not all the conceptual machinery.
Complex-systems and physical analogies are useful because they provide
organized ways to think about repeated interaction, cooperative locking,
threshold effects, and multiscale order. Their role in this chapter is
\textbf{medium} and explicitly interpretive.

\begin{remark}
The manuscript does not claim that a study desk, a research lab, and a neural
assembly are the same object. It claims that they may share analyzable
properties such as coupling, reinforcement, path dependence, and reorganized
cost landscapes.
\end{remark}

\subsection{Self-Organization and Order Parameters}

Classic complexity texts are useful here because they explain a foundational
idea: local interaction can produce higher-level coordination without requiring
a central executive that explicitly computes the full pattern in advance.

In the language of this manuscript, a framework can be treated as an
order-parameter-like summary of many local regularities:
\begin{itemize}[nosep]
  \item time of day,
  \item cue placement,
  \item posture,
  \item reward expectation,
  \item social synchronization,
  \item memory of prior success and prior failure.
\end{itemize}

Once these variables align often enough, the system no longer has to negotiate
the full behavioural choice from scratch. A macro-pattern exists, and that
pattern constrains the next round of micro-decisions.

\subsection{Path Dependence and Interference}

Path dependence is essential to the user requirement about ``past data
interference.'' The same present stimulus can lead to different action because
the system is carrying different traces of prior conditioning, prior losses,
prior rewards, or prior environmental pairings.

\begin{definition}[Path dependence]
\textbf{Path dependence} means that present transition costs depend on the
sequence of prior states and reinforcements, not only on the current snapshot.
\end{definition}

\begin{definition}[Interference]
\textbf{Interference} is the distortion or blocking of a target transition by
traces of previously reinforced routes, cues, or expectations that remain
active in the current context.
\end{definition}

\begin{discussion}
This is one of the chapter's most business-relevant points. Bad frameworks are
often not recreated because the present decision is morally wrong; they are
recreated because the old route still has more stored structure. Likewise, good
frameworks become easier not because the agent becomes ``better'' in the
abstract, but because the path has been rendered cheaper through repeated use.
\end{discussion}

\subsection{Multiscale Organization}

Emergence also matters because organization exists at several scales at once.
A useful framework can be:
\begin{enumerate}[nosep]
  \item locally enacted by small repeated actions,
  \item maintained by room layout, schedule structure, or social timing,
  \item embedded inside larger biological and conscious-state conditions.
\end{enumerate}

This multiscale viewpoint prevents a common mistake: trying to explain all
framework formation either by inner psychology alone or by environment alone.
The book's claim is more demanding. Frameworks form when several scales align
long enough to stabilize one route over others.
